% -*- mode:LaTex; mode:visual-line; mode:flyspell; fill-column:75-*-

\chapter{Mechatronics Design of Local-Sensing Swimming Robots} \label{chapdesign}

\section{Modular Undulatory Swimming Robot with Multi-sensory Feedback}
\subsection{Modular Structure Design}

\subsection{Electronics Design}

\subsection{Experimental Setup}


\subsection{Forward Swimming}

\subsection{C-Start Agile Escape Maneuver}







\section{A Robotics Platform for Drag Swimming}

Bioinspired swimming has gained considerable attention for efficient propulsion and maneuverability. However, complex system designs, unique actuation mechanisms, and modeling challenges have created barriers to entry for the broader robotics community. To address these issues, we introduce SwimGym: a low-cost platform for bioinspired swimming research. SwimGym comprises an untethered swimming robot, a portable hardware testbed, and a physics simulator based on the immersed boundary method. The robot design utilizes low-cost components and 3D printing to achieve a total cost less than \$200 USD. We demonstrate the robot's ability to achieve forward swimming with good sim-to-real agreement. The results highlight SwimGym as an effort towards a low-cost, open-source platform for bio-inspired swimming research that fully supports sim-to-real studies.
 
 
%%%%%%%%%%%%%%%%%%%%%%%%%%%%%%%%%%%%%%%%%%%%%%%%%%%%%%%%%%%%%%%%%%%%%%%%%%%%%%%%
 
\begin{figure}[t]
    \centering
    \includegraphics[width=\linewidth]{Figure/Figure1_11.pdf}
    \caption{An overview of SwimGym, a low-cost, open-source robotics platform for bio-inspired swimming research. Swimgym consists of an untethered cephalopod-inspired robot, a portable experimental testbed, and a simulation environment.}
    \label{fig:swimgym_overview}
\end{figure}

 
Aquatic animals have attracted considerable attention from biomechanics and robotics researchers due to their efficient and rapid locomotion in water~\cite{breder1926locomotion, Fish2014, van2019scaling, zhu2019tuna, Prakash_2024, wang2020development}. For example, the efficient, stable swimming of cephalopods and aquatic insects inspired various drag-based~\cite{kim2024development, kwak2023development, kwak2017design, hu2024design} mechanisms for robotic swimming. However, the diversity in actuator-mechanism design, specialized fabrication needs, and modeling difficulties pose significant challenges for reproducing these systems, which span different propulsion modes~\cite{hartmann2025highly, wu2024multimodal}. Furthermore, bio-inspired swimming involves complex fluid-robot interactions that are difficult to model for sim-to-real studies commonly found in robotics~\cite{angelidis2022gazebo, lee2023aquarium}.
 
Therefore, we propose \textbf{SwimGym}: a low-cost, open-source robot platform specifically designed for bio-inspired swimming research as shown in Figure\ref{fig:swimgym_overview}. SwimGym comprises an untethered swimming robot, a portable hardware testbed, and a corresponding simulation environment as shown in Figure~\ref{fig:swimgym_overview}. Utilizing COTS (commercial-off-the-shelf) components and 3D-printed parts, we design the robot from scratch to be modular and amenable to fabrication in house for under \$200 USD. In addition, a portable and reproducible testbed is provided for convenient testing of the robot in a fluid-filled environment with vision-based tracking. Finally, we utilize a unified-multiphysics approach to create a physics-based simulator of the test environment~\cite{lee2025unifiedframeworksimulatingstronglycoupled}. To showcase SwimGym's viability for swimming locomotion research, we demonstrate our robot achieving forward swimming in both simulation and on hardware. In future work, we plan to open-source the hardware design, experimental setup, and simulation code for use in downstream learning and controls tasks. In summary, SwimGym offers:
\begin{itemize}
    \item An untethered, insect-inspired swimming robot with a low build cost and entry-level fabrication process.
    \item A portable, low-cost experimental testbed for conducting real-world swimming locomotion.
    \item A physics-based, digital-twin simulator of the test environment for use in sim-to-real studies.
\end{itemize}
 
 
\section{Drag Based Swimming Platform}
 
\subsection{Robot Hardware and Experimental Testbed}
 
Inspired by aquatic insects and turtles---whose propulsion rely on drag forces generated by rowing strokes---the SwimGym robot features a main body with servo-actuated fins and rudders as shown in Figure~\ref{fig:hardware}. A gear assembly drives the fins symmetrically, while two miniature servos independently control each rudder via five-bar linkages for steering. The hardware assembly utilizes 3d-printed parts and commercially-available off-the-shelf (COTS) components as shown in Table~\ref{tab:parts_list}. The result is a low-cost (\$185) bill of material with an accessible fabrication process. An onboard, 7.4V battery powers the servos and onboard Adafruit Feather~M0 board, which integrates an ARM Cortex-M0 microcontroller with a low-power, long-range radio (LoRa) transceiver. This allows for truly untethered operation of the SwimGym robot while maintaining a stable, long-range wireless communication between the robot and a host computer.
 
To provide a portable swimming environment, we develop a low-cost experimental testbed consisting of a $122 \times 122 \times 40$\,cm pool and a vision-based tracking system as shown in Figure~\ref{fig:swimgym_overview}. The vision-based system requires only a (cellphone) camera and computer, which can also send motor commands wirelessly to the robot via LoRa.
 
\begin{figure*}[t]
    \centering
    \includegraphics[width=\textwidth]{Figure/Figure2_11.pdf}
    \caption{Hardware assembly of the SwimGym robot: a fully untethered swimmer with onboard power and microcontroller board. Servo motors drive fins for drag-based propulsion while rudders can be used to provide steering. Labels for each numbered part can be found in Table~\ref{tab:parts_list}.}
    \label{fig:hardware}
\end{figure*}
 
\begin{table}
\centering
\caption{SwimGym Robot Bill of Materials}
\label{tab:parts_list}
\begin{tabular}{c l c c}
\hline
\textbf{No.} & \textbf{Part Name} & \textbf{Quantity} & \textbf{Price (USD)} \\ \hline
1  & Main body for housing electronics         & 1        & \$10.80 \\
2  & Li-Ion 7.4V 2200mAh Battery               & 1        & \$8.50 \\
3  & PCB Board                                 & 1        & \$3.00 \\
4  & Mounting brackets for April tag           & 2        & \$0.31 \\
5  & Extra plate for pulsatile jet             & 1        & \$1.73 \\
6  & Gear Mounting Plate                       & 1        & \$0.39 \\
7--8 & Fins and Rudders                          & 2        & \$1.49 \\
9  & Gear assembly (0.5 module, 24 \& 48 teeth) & 1       & \$5.15 \\
10 & Transmission arms                         & 2        & \$3.20 \\
11 & M3 spacer (92871A009) and screw           & 1 set    & \$3.47 \\
12 & M10 20\,mm screw (head counterweight)     & 1        & \$0.78 \\
13 & Adafruit M0 controller (sending)          & 1        & \$39.95 \\
14 & Antennae for sending and receiving end     & 2        & \$1.90 \\
15 & Adafruit M0 controller (receiving)        & 1        & \$39.95 \\
16 & ES3059D servo motors                      & 2        & \$9.40 \\
17 & DC-to-DC Converters, OKR-T/6-W12-C         & 2        & \$19.06 \\
18 & Resistors (280\,$\Omega$, 320\,$\Omega$)   & 2        & \$0.25 \\
19 & ES3005HV servo motor                      & 1        & \$19.99 \\
20 & Screws, washers, helicoils, nuts, pins& Multiple & \$15.40 \\ 
\textbf{Total} & & &{\textbf{\$184.72}} \\
\hline
\end{tabular}
\end{table}
 
 
\subsection{Physics Simulator}
 
To facilitate sim-to-real studies, we develop a simulated gym environment replicating our testbed. To do so, we leverage variational mechanics to develop a Navier-Stokes-based fluid-structure interaction solver for robotics~\cite{lee2025unifiedframeworksimulatingstronglycoupled}. We reduce the sim-to-real gap by incorporating CAD-calculated mass properties into the simulation model.
 
To validate our simulator, we perform a sim-to-real demonstration of a forward-swimming gait. As shown in Figure~\ref{fig:simtoreal}, SwimGym's simulator is able to properly simulate the forward-swimming behavior consistent with the demonstration on hardware. Meanwhile, Genesis~\cite{genesis}, a state-of-the-art world simulator for robotics, is unable to simulate the achieved forward motion. Specifically, our framework achieves an RMSE error of 0.89~cm over the whole trajectory, a 90\% improvement over Genesis (8.86~cm RMSE error).
 
\begin{figure}[t]
    \centering
    \includegraphics[width=\columnwidth]{Figure/Figure6_02.pdf}
    \caption{Sim-to-real validation of the SwimGym robot executing an open-loop swimming gait in initially-still water. Our multiphysics simulator (\textcolor{green}{green}) is able to model the propulsion resulting from the fluid-robot interaction unlike the particle-based baseline in Genesis (black).}
    \label{fig:simtoreal}
\end{figure}